%----------------------------------------------------------------------------------------
% DE AFGELEIDE. Alle opgaven zijn eigen werk.
%
% Opbouw: Start, leerstof (theorie, werkwijze of uitleg, voorbeeld, korte oefening), en op een
% nieuwe pagina de Opdrachten (door elkaar, stijgende moeilijkheid), Koppel en Test jezelf.
% Labels van dit hoofdstuk beginnen met h2-.
%----------------------------------------------------------------------------------------

\mychapter{1}{De afgeleide}

\begin{probleem}
\naast{Je schenkt een cocktailglas vol, gelijkmatig $10$ ml per seconde. Het glas is een kegel met de punt onderaan: onderaan is het smal, dus daar stijgt het vocht snel, en hogerop steeds trager. Uit de vorm van het glas volgt dat het vocht na $t$ seconden
\[h(t) = 3\sqrt[3]{t} \quad \text{cm}\]
hoog staat, tot het glas na ongeveer $19$ seconden vol is. Na $1$ s staat het vocht $3$ cm hoog, na $8$ s $6$ cm: in de eerste seconde stijgt het $3$ cm, in de zeven seconden daarna samen ook $3$ cm.

Hoe snel stijgt het vocht op het ogenblik $t = 8$, in cm per seconde? En wat betekent \textit{hoe snel} eigenlijk, op één ogenblik?}{%
\begin{tikzpicture}[scale=.36,font=\scriptsize]
  \fill[accent!25] (0,0) -- (-3.57,6) -- (3.57,6) -- cycle;
  \draw[accent!70!black] (-3.57,6) -- (3.57,6);
  \draw[grijs,thick] (-4.76,8) -- (0,0) -- (4.76,8);
  \draw[grijs,thick] (0,0) -- (0,-3.2);
  \draw[grijs,thick] (-2.4,-3.2) -- (2.4,-3.2);
  \draw[grijs,<->] (5.6,0) -- (5.6,6) node[midway,right,text=black] {$h$};
  \draw[grijs,densely dotted] (0,0) -- (5.9,0) (3.57,6) -- (5.9,6);
\end{tikzpicture}}
\end{probleem}

\begin{alleenvoor}{2,3}
\lkr{Leerplan \pakketnaam: LPD \alleenvoor{2}{18 (de afgeleide als limiet van een differentiequotiënt en als richtingscoëfficiënt van de raaklijn), LPD 19 (het grafische verband tussen $f$ en $f'$) en het begin van LPD 20 (de regels voor machten, som, product en quotiënt)}\alleenvoor{3}{19 (de afgeleide als limiet van een differentiequotiënt en als richtingscoëfficiënt van de raaklijn), LPD 20 (het grafische verband tussen $f$ en $f'$) en het begin van LPD 21 (de regels voor machten, som, product en quotiënt)}. Een intuïtief limietbegrip volstaat (wenk bij LPD \alleenvoor{2}{18}\alleenvoor{3}{19}): het hoofdstuk Limieten en continuïteit valt in dit pakket weg, en leerlingen benaderen de afgeleide ook met een tabel (ICT). De flitspaal en de trajectcontrole, de contexten (debiet, marginale kost\alleenvoor{2}{, elasticiteit}, helling) en het schetsen van de grafiek van $f'$ komen uit de wenken.\alleenvoor{2}{ In B+S is bewijsvoering bij LPD 4 een extra: de rekenregels worden aannemelijk gemaakt, niet bewezen. De afgeleide van $x$ en $x^2$ met de definitie (het voorbeeld dat LPD 4 noemt) blijft.}\alleenvoor{3}{ In B+S' hoort bewijsvoering bij LPD 4: de bewijzen van de rekenregels blijven.} Het beeld van de voorkant (koorden door $P$ die naar de raaklijn draaien) is de figuur in sectie 1. Laat leerlingen de definitie een paar keer zelf uitwerken voor ze de rekenregels gebruiken. Het glas in het probleem bevat $\frac{10}{27}h^3$ ml als het vocht $h$ cm hoog staat (een kegel met straal $0{,}59\,h$ op hoogte $h$); bij een diepte van $8$ cm is het bovenaan $9{,}5$ cm breed en bevat het $190$ ml. Uit $10t = \frac{10}{27}h^3$ volgt $h = 3\sqrt[3]{t}$. De oplossing staat in voorbeeld \ref{vb:h2-terug}, op het einde van de leerstof.}
\end{alleenvoor}
\begin{alleenvoor}{1}
\lkr{Leerplan Basis: LPD 14 (de afgeleide als limiet van een differentiequotiënt en als richtingscoëfficiënt van de raaklijn) en het begin van LPD 15 (van de afgeleide in een punt naar de afgeleide functie, grafisch). Een intuïtief limietbegrip volstaat: leerlingen benaderen de afgeleide met een tabel van differentiequotiënten (ICT) en rekenen ze alleen exact uit bij eerste- en tweedegraadsfuncties, de uitbreiding bij LPD 15. De rekenregels, de regel voor machten en het verband tussen afleidbaar en continu vallen buiten Basis. Het beeld van de voorkant (koorden door $P$ die naar de raaklijn draaien) is de figuur in sectie 1. Het glas in het probleem bevat $\frac{10}{27}h^3$ ml als het vocht $h$ cm hoog staat (een kegel met straal $0{,}59\,h$ op hoogte $h$); uit $10t = \frac{10}{27}h^3$ volgt $h = 3\sqrt[3]{t}$. De oplossing staat in voorbeeld \ref{vb:h2-terug}, op het einde van sectie 1: met een tabel, en met een controle uit de vorm van het glas.}
\end{alleenvoor}

Een snelheid is een verandering gedeeld door een tijd. Over een tijdsduur kun je ze uitrekenen, maar op één ogenblik verloopt er geen tijd, en stijgt het vocht in het glas ook niet. Toch toont een snelheidsmeter hoe snel je \textit{nu} rijdt, niet gemiddeld over de rit. Hoe kan dat? Het antwoord is een limiet: neem de gemiddelde snelheid over een steeds kortere tijd. Zo ontstaat de afgeleide, misschien het belangrijkste begrip van de analyse. \alleenvoor{2,3}{Met de rekenregels op het einde van het hoofdstuk lossen we het probleem van het glas op.}\alleenvoor{1}{Op het einde van sectie 1 lossen we het probleem van het glas op, met een tabel die je rekentoestel of GeoGebra maakt.}

\begin{oefeningen}[Start]
\oef{h2-start-1}

\oef{h2-verken}
\end{oefeningen}

%========================================================================================
\section{De afgeleide in een punt}

\subsection{Het differentiequotiënt}

Begin bij het glas uit het probleem. Tussen $t = 1$ en $t = 8$ stijgt het vocht van $3$ naar $6$ cm, gemiddeld $\frac{6 - 3}{8 - 1} = \frac37 \approx 0{,}43$ cm per seconde. Dat is een gemiddelde over zeven seconden, en in die tijd vertraagt het vocht sterk. Neem daarom een steeds korter interval vanaf $t = 8$. Tussen $8$ en $9$ s is de gemiddelde snelheid $\frac{h(9) - h(8)}{9 - 8} \approx 0{,}240$ cm/s. Met kortere intervallen vind je:

\figuren{{\renewcommand{\arraystretch}{1.2}%
\begin{tabular}{c|ccc}
van $t = 8$ tot & $9$ & $8{,}1$ & $8{,}01$ \\ \hline
gemiddelde snelheid (cm/s) & $0{,}240$ & $0{,}249$ & $0{,}2499$
\end{tabular}}}

De gemiddelde snelheden lijken naar $0{,}25$ cm/s te gaan. Een interval met lengte $0$ kan niet: dan deel je $0$ door $0$. Toch zie je in de tabel een getal verschijnen: hoe korter het interval, hoe dichter de gemiddelde snelheid bij $0{,}25$ komt. We zeggen dat $0{,}25$ de \textit{limiet} is. De snelheid op één ogenblik is zo'n limiet. Op de grafiek van $h$ is elke gemiddelde snelheid de richtingscoëfficiënt van een rechte door twee punten. Bekijk dat nu voor een willekeurige functie.

\naast{%
Neem twee punten $P(a, f(a))$ en $Q(a + h, f(a + h))$ op de grafiek van $f$. De rechte erdoor is een \textbf{koorde}. Haar richtingscoëfficiënt is de verandering van $f$ gedeeld door de verandering van $x$: hoeveel $f$ gemiddeld verandert per eenheid, tussen $a$ en $a + h$.

Laat $Q$ naar $P$ schuiven ($h \to 0$). De koorden draaien dan naar een grensstand: de \textbf{raaklijn} in $P$.}{%
\begin{tikzpicture}[x=1.1cm,y=.55cm,font=\scriptsize]
  \draw[->,lijngrijs] (-.3,0) -- (2.6,0) node[right,text=black] {$x$};
  \draw[->,lijngrijs] (0,-.3) -- (0,5.6) node[above,text=black] {$y$};
  \kromme{0}{2.3}{(\x)^2}
  \draw[grijs] (.55,-.35) -- (2.3,4.9);
  \draw[grijs,dashed] (.5,-.25) -- (2.3,4.25);
  \draw[accent,thick] (.3,-.4) -- (2.2,3.4);
  \fill[accent] (1,1) circle (.05) node[above left,text=black] {$P$};
  \fill (2,4) circle (.05) node[right] {$Q$};
  \fill (1.5,2.25) circle (.04);
  \node[accent,right] at (2.2,3.3) {raaklijn};
\end{tikzpicture}}

\begin{theorie}
Het \textbf{\textcolor{accent}{differentiequotiënt}} van $f$ tussen $a$ en $a + h$ (met $h \neq 0$) is
\[\formule{\frac{\Delta y}{\Delta x} = \frac{f(a + h) - f(a)}{h}}.\]
Het is de richtingscoëfficiënt van de koorde door $(a, f(a))$ en $(a + h, f(a + h))$.
\end{theorie}

De letter $\Delta$ (delta) staat voor \textit{verschil}: $\Delta x = h$ is de verandering van $x$, en $\Delta y = f(a + h) - f(a)$ de verandering van de functiewaarde. Het getal $h$ mag ook negatief zijn: dan ligt $Q$ links van $P$.

\voorbeeld{}{Bereken het differentiequotiënt van $f(x) = x^2$ tussen $1$ en $1 + h$. Wat is het voor $h = 2$, $h = 1$ en $h = 0{,}1$?\label{vb:h2-dq}}{%
\uitleg{Werk het differentiequotiënt eerst algemeen uit, met de letter $h$. Daarna vul je de getallen in, en zie je wat er gebeurt als $h$ kleiner wordt.}
\begin{stappen}
\frac{f(1 + h) - f(1)}{h} = \frac{(1 + h)^2 - 1}{h} & De definitie, met $a = 1$ en $f(1) = 1$. \\[1ex]
= \frac{2h + h^2}{h} = 2 + h & Werk het kwadraat uit en deel door $h$; dat mag, want $h \neq 0$. \\[1ex]
{4, \quad 3, \quad 2{,}1} & Invullen van $h = 2$, $h = 1$ en $h = 0{,}1$ geeft de richtingscoëfficiënten van drie koorden door $P(1, 1)$. Voor $h = 2$ is dat de koorde naar $Q(3, 9)$.
\end{stappen}
Hoe kleiner $h$, hoe dichter het differentiequotiënt bij $2$ komt. De koorden naderen een rechte met richtingscoëfficiënt $2$.}

\begin{oefeningen}
\oef{h2-dq}
\end{oefeningen}

\subsection{De afgeleide in een punt}

Wat de koorden naderen, is een limiet van differentiequotiënten. Die limiet krijgt een eigen naam.

\begin{theorie}
De \textbf{\textcolor{accent}{afgeleide}} van $f$ in $a$ is
\[\formule{f'(a) = \lim_{h \to 0} \frac{f(a + h) - f(a)}{h}},\]
als die limiet bestaat. De functie $f$ is dan \textbf{\textcolor{accent}{afleidbaar}} in $a$.
\end{theorie}

De notatie $\lim_{h \to 0}$ lees je als \textit{de limiet voor $h$ naar $0$}: het getal waar het differentiequotiënt naartoe gaat als $h$ steeds dichter bij $0$ komt, langs beide kanten, zonder ooit $0$ te worden. Invullen van $h = 0$ kan dus niet: dat geeft $\frac00$. Toch vind je de limiet, op twee manieren. Met ICT maak je een tabel van het differentiequotiënt voor steeds kleinere $h$, en kijk je waar de uitkomsten naartoe gaan. Bij een eenvoudige functie kan het ook exact.

\voorbeeld{}{Benader $f'(1)$ voor $f(x) = 2^x$ met een tabel. Geef ook de raaklijn in het punt $(1, 2)$.\label{vb:h2-tabel}}{%
\naast{%
\uitleg{Laat het rekentoestel het differentiequotiënt uitrekenen voor kleine $h$, positief en negatief. In de tabel van de TI-84 speelt \texttt{X} de rol van $h$.}
\begin{stappenvrij}{3.6cm}{5.6cm}
\frac{2^{1 + h} - 2}{h} & Het differentiequotiënt in $a = 1$, met $f(1) = 2$. Voer in: \texttt{Y1 = (2\textasciicircum(1+X)-2)/X}, en zet \texttt{TBLSET} op \texttt{Ask}.
\end{stappenvrij}
\vspace{1mm}
{\renewcommand{\arraystretch}{1.2}\small%
\begin{tabular}{c|cccccc}
$h$ & $0{,}1$ & $0{,}01$ & $0{,}001$ & $-0{,}001$ & $-0{,}01$ & $-0{,}1$ \\ \hline
& $1{,}4355$ & $1{,}3911$ & $1{,}3868$ & $1{,}3858$ & $1{,}3815$ & $1{,}3393$
\end{tabular}}
\vspace{1mm}
\begin{stappenvrij}{3.6cm}{5.6cm}
f'(1) \approx 1{,}386 & Langs beide kanten gaan de uitkomsten naar hetzelfde getal. Controle: \texttt{2nd} \texttt{CALC} \texttt{6:dy/dx} in $x = 1$, of \texttt{f'(1)} in GeoGebra, geeft $1{,}386\,29$. \\
y = 2 + 1{,}386\,(x - 1) & De raaklijn: de rechte door $(1, 2)$ met die richtingscoëfficiënt.
\end{stappenvrij}}{%
\begin{tikzpicture}[x=.75cm,y=.55cm,font=\scriptsize]
  \rooster{-1}{1}{2}{-1}{1}{5}
  \kromme{-1}{2.3}{2^(\x)}
  \draw[grijs,thick] (-.4,.06) -- (2.4,3.94);
  \vol{(1,2)}
\end{tikzpicture}}}

\begin{werkwijze}
Voer het differentiequotiënt $\frac{f(a + h) - f(a)}{h}$ in, met $h$ als veranderlijke. Bekijk de uitkomsten voor $h = \pm 0{,}1$, $\pm 0{,}01$ en $\pm 0{,}001$. Gaan ze langs beide kanten naar hetzelfde getal, dan is dat bij benadering $f'(a)$. Controleer met \texttt{dy/dx} op de TI-84 of met \texttt{f'(a)} in GeoGebra.
\end{werkwijze}

\begin{oefeningen}
\oef{h2-tabel}
\end{oefeningen}

Bij \alleenvoor{1}{een tweedegraadsfunctie}\alleenvoor{2,3}{veelterm- en rationale functies} kun je de limiet ook exact berekenen. Vereenvoudig dan eerst het differentiequotiënt, tot de $h$ in de noemer wegvalt. Daarna mag je $h = 0$ invullen.

\voorbeeld{}{Bereken met de definitie $f'(3)$ voor $f(x) = x^2$.\label{vb:h2-definitie}}{%
\uitleg{Werk het differentiequotiënt uit tot de $h$ in de noemer wegvalt. Neem dan de limiet.}
\begin{stappen}
\frac{f(3 + h) - f(3)}{h} = \frac{(3 + h)^2 - 9}{h} & De definitie, met $f(3) = 9$. \\[1ex]
= \frac{9 + 6h + h^2 - 9}{h} = \frac{6h + h^2}{h} & Werk het kwadraat uit. De $9$ valt weg, zoals altijd: $f(a)$ verdwijnt in de teller. \\[1ex]
= 6 + h & Deel door $h$; dat mag, want $h \neq 0$. \\[1ex]
f'(3) = \lim_{h \to 0} (6 + h) = 6 & Nu kunnen we invullen.
\end{stappen}}

\begin{werkwijze}
Bereken $f(a + h)$ en $f(a)$, en werk de teller $f(a + h) - f(a)$ uit.\alleenvoor{2,3}{ Breng zo nodig op één noemer\alleenvoor{3}{, of vermenigvuldig met de toegevoegde}.} Zonder een factor $h$ af en schrap die tegen de noemer. Neem dan de limiet voor $h \to 0$ door in te vullen.
\end{werkwijze}

\begin{oefeningen}
\oef{h2-afgpunt}
\end{oefeningen}

\subsection{De raaklijn}

De koorden door $P$ draaien naar een grensstand. Die rechte gaat door $P$ en heeft als richtingscoëfficiënt de limiet van de richtingscoëfficiënten van de koorden: de afgeleide.

\begin{theorie}
Is $f$ afleidbaar in $a$, dan is de \textbf{\textcolor{accent}{raaklijn}} aan de grafiek van $f$ in $(a, f(a))$ de rechte door dat punt met richtingscoëfficiënt $f'(a)$:
\[\formule{y = f(a) + f'(a)\,(x - a)}.\]
\end{theorie}

Dat is de vergelijking van een rechte door een punt met een gegeven richtingscoëfficiënt, die je al kent. Nieuw is alleen dat de richtingscoëfficiënt een afgeleide is.

\voorbeeld{}{Geef de raaklijn aan de grafiek van $f(x) = x^2$ in het punt $(3, 9)$.\label{vb:h2-raaklijn}}{%
\naast{%
\uitleg{De richtingscoëfficiënt vonden we in het vorige voorbeeld: $f'(3) = 6$. Vul in de vergelijking van de raaklijn in.}
\begin{stappenvrij}{4.3cm}{5cm}
y = f(3) + f'(3)\,(x - 3) & De vergelijking, met $a = 3$. \\
y = 9 + 6(x - 3) & $f(3) = 9$ en $f'(3) = 6$. \\
y = 6x - 9 & Uitwerken.
\end{stappenvrij}
In de figuur zie je dat de rechte de parabool in $(3, 9)$ raakt, zonder haar te snijden.}{%
\begin{tikzpicture}[x=.6cm,y=.2cm,font=\scriptsize]
  \draw[black!14,very thin,xstep=1,ystep=2] (0,-4) grid (4,16);
  \draw[->,grijs] (0,0) -- (4.4,0) node[right,text=black] {$x$};
  \draw[->,grijs] (0,-4) -- (0,16.8) node[above,text=black] {$y$};
  \foreach \t in {1,2,3,4}{\node[below,inner sep=1.5pt] at (\t,0) {$\t$};}
  \foreach \t in {4,8,12,16}{\node[left,inner sep=1.5pt] at (0,\t) {$\t$};}
  \kromme{0}{4}{(\x)^2}
  \draw[grijs,thick] (.84,-4) -- (4,15);
  \vol{(3,9)}
\end{tikzpicture}}}

\begin{werkwijze}
Bereken $f(a)$ en $f'(a)$. Vul ze in $y = f(a) + f'(a)(x - a)$ in en werk uit tot $y = mx + q$.
\end{werkwijze}

\begin{oefeningen}
\oef{h2-defpunt}
\end{oefeningen}

\subsection{Ogenblikkelijke snelheid}

Het differentiequotiënt heeft een betekenis in elke context. Is $s(t)$ de afgelegde weg op tijdstip $t$, dan is $s(t + h) - s(t)$ de weg die je aflegt tussen $t$ en $t + h$. Gedeeld door de tijd $h$ is dat de gemiddelde snelheid.

\begin{theorie}
Is $s(t)$ de afgelegde weg op tijdstip $t$, dan is de \textbf{\textcolor{accent}{gemiddelde snelheid}} tussen $t$ en $t + h$ gelijk aan $\formule{\frac{s(t + h) - s(t)}{h}}$, en de \textbf{\textcolor{accent}{ogenblikkelijke snelheid}} op tijdstip $t$ gelijk aan $\formule{v(t) = s'(t)}$.
\end{theorie}

Zo lost de afgeleide de vraag uit de inleiding op: de snelheid op één ogenblik is de limiet van de gemiddelde snelheid over steeds kortere tijd. Voor het glas is de stijgsnelheid op het ogenblik $t$ dus $h'(t)$. Met de definitie is die moeilijk\alleenvoor{1}{ exact} te berekenen, want $h$ bevat een derdemachtswortel. \alleenvoor{2,3}{Met een tabel kun je haar benaderen, maar exact lukt het pas met de regels uit de volgende secties.}\alleenvoor{1}{Met een tabel lukt het wel: dat doen we op het einde van deze sectie.}

Zo meten een flitspaal en een trajectcontrole iets anders. Een trajectcontrole deelt de afstand tussen twee camera's door de tijd die je erover doet: een gemiddelde snelheid, een differentiequotiënt. Een flitspaal meet je snelheid op één ogenblik: een afgeleide. Wie op een traject gemiddeld $70$ km/h rijdt, kan onderweg dus even $90$ km/h gereden hebben. Op de grafiek van de afgelegde weg is de eerste de richtingscoëfficiënt van een koorde, de tweede die van een raaklijn.

\voorbeeld{}{Een steen valt; na $t$ seconden heeft hij $s(t) = 5t^2$ meter afgelegd. Hoe snel valt hij na 2 seconden?\label{vb:h2-snelheid}}{%
\naast{%
\uitleg{Neem de gemiddelde snelheid tussen $2$ en $2 + h$, en laat $h$ naar $0$ gaan.}
\begin{stappenvrij}{5.5cm}{4.6cm}
{\frac{s(2 + h) - s(2)}{h} = \frac{5(2 + h)^2 - 20}{h}} & De gemiddelde snelheid, met $s(2) = 20$. \\[1ex]
{= \frac{20h + 5h^2}{h} = 20 + 5h} & Uitwerken en delen door $h$. \\[1ex]
{s'(2) = \lim_{h \to 0} (20 + 5h) = 20} & In m/s, ongeveer 72 km/h.
\end{stappenvrij}
In de figuur is $20$ de richtingscoëfficiënt van de raaklijn aan de grafiek van $s$ in $t = 2$.}{%
\begin{tikzpicture}[x=.75cm,y=.06cm,font=\scriptsize]
  \draw[black!14,very thin,xstep=1,ystep=10] (0,0) grid (3,45);
  \draw[->,grijs] (0,0) -- (3.4,0) node[right,text=black] {$t$};
  \draw[->,grijs] (0,0) -- (0,48) node[above,text=black] {$s$};
  \foreach \t in {1,2,3}{\node[below,inner sep=1.5pt] at (\t,0) {$\t$};}
  \foreach \t in {10,20,30,40}{\node[left,inner sep=1.5pt] at (0,\t) {$\t$};}
  \kromme{0}{3}{5*(\x)^2}
  \draw[grijs,thick] (1,0) -- (3.2,44);
  \vol{(2,20)}
\end{tikzpicture}}}

\begin{oefeningen}
\oef{h2-definitie}
\end{oefeningen}

\subsection{De afgeleide in een context}

Niet alleen een snelheid is een afgeleide. Telkens als een grootheid verandert met een andere, zegt de afgeleide hoe snel dat gebeurt, op één ogenblik of bij één waarde.

\begin{theorie}
Is $y = f(x)$, dan is $\formule{f'(a)}$ de \textbf{\textcolor{accent}{ogenblikkelijke verandering}} van $y$ bij $x = a$. Haar eenheid is de $\formule{\text{eenheid van } y \text{ per eenheid van } x}$.
\end{theorie}

Zo is de afgeleide van een volume naar de tijd een \textit{debiet}, in liter per minuut. De afgeleide van de kosten naar het aantal stuks is de \textit{marginale kost}, in euro per stuk: ongeveer wat één stuk extra kost. En de afgeleide van de hoogte naar de afgelegde afstand is de \textit{helling} van een weg.

\voorbeeld{}{Tijdens een onweer stroomt er water in een stuwmeer. $V(t)$ is het volume in miljoen m$^3$, $t$ uur na het begin van de bui. Er geldt $V(6) = 40$ en $V'(6) = 0{,}3$. Wat betekent $V'(6)$? Schat het volume een halfuur later.\label{vb:h2-context}}{%
\uitleg{Begin bij de eenheid: die zegt welke verandering de afgeleide meet.}
\begin{stappen}
\text{miljoen m}^3 \text{ per uur} & De eenheid van $V$ gedeeld door die van $t$. \\[1ex]
V'(6) = 0{,}3 & Zes uur na het begin van de bui stroomt er netto $0{,}3$ miljoen m$^3$ per uur bij: het debiet. Het teken is positief: het meer vult zich. \\[1ex]
V(6{,}5) \approx 40 + 0{,}3 \cdot 0{,}5 = 40{,}15 & Blijft het debiet ongeveer gelijk, dan komt er in een halfuur ongeveer $0{,}15$ miljoen m$^3$ bij.
\end{stappen}
De schatting is maar goed zolang het debiet weinig verandert: ze volgt de raaklijn, niet de grafiek.}

\begin{oefeningen}
\oef{h2-context}
\end{oefeningen}
\begin{alleenvoor}{2}

In de economie meet men een verandering vaak in procent. Stijgt de prijs $p$ van een product met $1\,\%$, met hoeveel procent verandert dan de gevraagde hoeveelheid $q(p)$? Dat getal heet de \textit{elasticiteit} van de vraag. Een prijsstijging van $1\,\%$ is $\Delta p = 0{,}01\,p$, en daardoor verandert de vraag ongeveer met $\Delta q \approx q'(p) \cdot 0{,}01\,p$, zoals bij de schatting met de raaklijn hierboven. Gedeeld door $q(p)$ is dat $\frac{\Delta q}{q(p)} \approx 0{,}01 \cdot \frac{p}{q(p)}\,q'(p)$: de vraag verandert met ongeveer $\frac{p}{q(p)}\,q'(p)$ procent.

\begin{theorie}
De \textbf{\textcolor{accent}{elasticiteit}} van de vraag $q(p)$ bij de prijs $p$ is \quad $\formule{E = \frac{p}{q(p)} \cdot q'(p)}$.
\end{theorie}

Meestal is $E$ negatief: wie de prijs verhoogt, verkoopt minder. Is $|E| > 1$, dan heet de vraag \textit{elastisch}: de vraag daalt procentueel sterker dan de prijs stijgt.

\voorbeeld{}{De vraag naar een concertticket is $q(p) = 2000 - 25p$ (in stuks, $p$ in euro). Bereken de elasticiteit bij een prijs van $30$ euro. Wat betekent ze?\label{vb:h2-elasticiteit}}{%
\uitleg{Bereken $q(30)$ en de afgeleide $q'(30)$, en vul in. De grafiek van $q$ is een rechte: de afgeleide is haar richtingscoëfficiënt.}
\begin{stappen}
q(30) = 2000 - 750 = 1250, \quad q'(30) = -25 & Bij $30$ euro verkoop je $1250$ tickets; per euro meer zijn het er $25$ minder. \\[1ex]
E = \frac{30}{1250} \cdot (-25) = -0{,}6 & De elasticiteit.
\end{stappen}
Stijgt de prijs met $1\,\%$, dan daalt de vraag met ongeveer $0{,}6\,\%$: de vraag is niet elastisch. De ontvangst $p \cdot q(p)$ stijgt dus nog als de organisator de prijs wat verhoogt.}

\begin{oefeningen}
\oef{h2-elasticiteit}
\end{oefeningen}
\end{alleenvoor}
\begin{alleenvoor}{1}

\needspace{22\baselineskip}
\subsection{Terug naar het probleem}

Nu kunnen we de stijgsnelheid in het glas benaderen, met een tabel zoals bij $2^x$.

\voorbeeld{}{Terug naar het probleem. Je schenkt een kegelvormig cocktailglas vol, $10$ ml per seconde. Na $t$ seconden staat het vocht $h(t) = 3\sqrt[3]{t}$ cm hoog. Hoe snel stijgt het vocht na $8$ seconden? En na $1$ seconde?\label{vb:h2-terug}}{%
\naast{%
\uitleg{De stijgsnelheid op een ogenblik is de afgeleide $h'(t)$. Noem de kleine tijdstap $\Delta t$, want de letter $h$ is hier de hoogte. Maak een tabel voor kleine $\Delta t$, langs beide kanten.}
\begin{stappenvrij}{3.8cm}{4.6cm}
\frac{h(8 + \Delta t) - 6}{\Delta t} & Het differentiequotiënt in $8$, met $h(8) = 3 \cdot 2 = 6$. Voer in: \texttt{Y1 = (3(8+X)\textasciicircum(1/3)-6)/X}. \\
0{,}249\,99 \text{ en } 0{,}250\,01 & De uitkomsten voor $\Delta t = 0{,}001$ en $\Delta t = -0{,}001$. \\
h'(8) = 0{,}25 & Na $8$ s stijgt het vocht $0{,}25$ cm per seconde. \\
h'(1) = 1 & Dezelfde tabel in $t = 1$ geeft $0{,}999\,7$ en $1{,}000\,3$: na $1$ s stijgt het vocht $1$ cm per seconde, vier keer zo snel.
\end{stappenvrij}}{%
\begin{tikzpicture}[x=.25cm,y=.45cm,font=\scriptsize]
  \draw[black!14,very thin,xstep=1cm,ystep=.9cm] (0,0) grid (20,8);
  \draw[->,grijs] (0,0) -- (21,0) node[right,text=black] {$t$};
  \draw[->,grijs] (0,0) -- (0,8.8) node[above,text=black] {$h$};
  \foreach \t in {4,8,12,16}{\node[below,inner sep=1.5pt] at (\t,0) {$\t$};}
  \foreach \t in {2,4,6,8}{\node[left,inner sep=1.5pt] at (0,\t) {$\t$};}
  \draw[accent,thick] plot[domain=0:8,samples=60,smooth] ({(\x/3)^3},\x);
  \draw[grijs,thick] (0,4) -- (16,8);
  \vol{(8,6)}
\end{tikzpicture}}
In de figuur is $\frac14$ de richtingscoëfficiënt van de raaklijn in $(8, 6)$. Dat klopt met de tabel aan het begin van deze sectie, en ook met de vorm van het glas. Op $6$ cm hoogte is het vochtoppervlak $40$ cm$^2$, en $10$ cm$^3$ per seconde verdeeld over $40$ cm$^2$ geeft een laagje van $0{,}25$ cm per seconde. Onderaan is het oppervlak klein, daarom stijgt het vocht daar zo snel. In $t = 0$ vertrekt de grafiek zelfs verticaal: daar bestaat de afgeleide niet.}
\end{alleenvoor}

%========================================================================================
\section{De afgeleide functie}

\subsection{Van getal naar functie}

Tot nu toe berekenden we de afgeleide in één punt. Doe je dat in elk punt $x$, dan hoort bij elke $x$ een getal $f'(x)$: een nieuwe functie. Ze zegt in elk punt hoe steil de grafiek van $f$ is.

\begin{theorie}
De \textbf{\textcolor{accent}{afgeleide functie}} van $f$ is
\[\formule{f'(x) = \lim_{h \to 0} \frac{f(x + h) - f(x)}{h}}.\]
Andere notaties: $\formule{\frac{\mathrm{d}f}{\mathrm{d}x}}$ of, met $y = f(x)$, $\formule{\frac{\mathrm{d}y}{\mathrm{d}x}}$.
\end{theorie}

De notatie $\frac{\mathrm{d}y}{\mathrm{d}x}$ is van Leibniz. Ze herinnert aan het differentiequotiënt $\frac{\Delta y}{\Delta x}$: een heel kleine verandering van $y$ gedeeld door een heel kleine verandering van $x$. Het is geen echte breuk\alleenvoor{2,3}{, maar ze rekent soms wel zo, zoals je in hoofdstuk 2 ziet}.

\naast{%
In de figuur stijgt $f$ waar $f'$ positief is, en daalt $f$ waar $f'$ negatief is. In de toppen is de raaklijn horizontaal, en daar is $f'$ nul. Hoe steiler de grafiek van $f$, hoe verder die van $f'$ van de $x$-as ligt.}{%
\begin{tikzpicture}[x=.55cm,y=.42cm,font=\scriptsize]
  \rooster{-2}{1}{2}{-2}{1}{3}
  \draw[grijs,dashed] (-1,-2) -- (-1,3) (1,-2) -- (1,3);
  \kromme{-2.2}{2.2}{(\x)^3/3-\x}
  \node[accent,above] at (-1.3,1) {$f$};
\end{tikzpicture}\quad
\begin{tikzpicture}[x=.55cm,y=.42cm,font=\scriptsize]
  \rooster{-2}{1}{2}{-2}{1}{3}
  \draw[grijs,dashed] (-1,-2) -- (-1,3) (1,-2) -- (1,3);
  \kromme[grijs]{-2.05}{2.05}{(\x)^2-1}
  \node[grijs,right] at (-1.6,2.9) {$f'$};
\end{tikzpicture}}

\voorbeeld{}{Bereken met de definitie de afgeleide functie van $f(x) = x^2$.\label{vb:h2-afgfunctie}}{%
\uitleg{Dezelfde berekening als in een punt, maar met de letter $x$ in plaats van een getal.}
\begin{stappen}
\frac{f(x + h) - f(x)}{h} = \frac{(x + h)^2 - x^2}{h} & De definitie. \\[1ex]
= \frac{2xh + h^2}{h} = 2x + h & Uitwerken, en delen door $h \neq 0$. \\[1ex]
f'(x) = \lim_{h \to 0} (2x + h) = 2x & Voor $h \to 0$ blijft $x$ gewoon staan.
\end{stappen}
Zo hoeven we niet meer voor elk punt opnieuw te rekenen: $f'(3) = 6$ uit het vorige voorbeeld volgt nu meteen.}

\begin{oefeningen}
\oef{h2-afgfunctie}
\end{oefeningen}

\begin{alleenvoor}{1}
\subsection{Eerste- en tweedegraadsfuncties}

Bij $x^2$ vonden we $2x$. Moet je die berekening voor elke tweedegraadsfunctie opnieuw maken? Nee: ze verloopt altijd op dezelfde manier, en geeft een eenvoudige formule.

\begin{theorie}
\begin{stelling}{Afgeleide van een eerste- en een tweedegraadsfunctie}
$\formule{(ax + b)' = a}$ \qquad en \qquad $\formule{(ax^2 + bx + c)' = 2ax + b}$.
\end{stelling}
\begin{bewijs}
Neem $f(x) = ax^2 + bx + c$. Dan is
\bewijsregel{f(x + h) - f(x) = a(x^2 + 2xh + h^2) + b(x + h) + c - ax^2 - bx - c = 2axh + ah^2 + bh ,}
dus het differentiequotiënt is $2ax + b + ah$ voor $h \neq 0$, met limiet $2ax + b$. Voor $a = 0$ staat er de eerstegraadsfunctie $bx + c$, met afgeleide $b$: dat is de eerste formule, met andere letters.
\end{bewijs}
\end{theorie}

De grafiek van een eerstegraadsfunctie is een rechte: overal even steil, met de richtingscoëfficiënt als helling. De afgeleide van een tweedegraadsfunctie is een eerstegraadsfunctie. Ze is nul in één punt: daar is de raaklijn horizontaal, in de top van de parabool.

\voorbeeld{}{Bereken de afgeleide van $f(x) = 3x^2 - 12x + 5$. In welk punt is de raaklijn aan de grafiek horizontaal?\label{vb:h2-tweedegraad}}{%
\naast{%
\uitleg{Pas de formule toe met $a = 3$, $b = -12$ en $c = 5$. Een horizontale raaklijn heeft richtingscoëfficiënt $0$.}
\begin{stappenvrij}{4.8cm}{5.1cm}
f'(x) = 2 \cdot 3x - 12 = 6x - 12 & De formule: $2ax + b$. De constante $5$ verdwijnt: ze verschuift de grafiek alleen. \\
6x - 12 = 0 \iff x = 2 & Daar is de raaklijn horizontaal. \\
f(2) = 12 - 24 + 5 = -7 & In $(2, -7)$, de top van de dalparabool.
\end{stappenvrij}
Links van $2$ is $f'(x) < 0$ en daalt de grafiek, rechts ervan is $f'(x) > 0$ en stijgt ze.}{%
\begin{tikzpicture}[x=.55cm,y=.2cm,font=\scriptsize]
  \rooster{-1}{1}{5}{-8}{2}{10}
  \kromme{-.35}{4.35}{3*(\x)^2-12*\x+5}
  \draw[grijs,thick] (.6,-7) -- (3.4,-7);
  \vol{(2,-7)}
\end{tikzpicture}}}

\begin{oefeningen}
\oef{h2-tweedegraad}
\end{oefeningen}
\end{alleenvoor}

\subsection{De grafiek van \texorpdfstring{$f'$}{f'} schetsen}

Van veel functies ken je geen voorschrift: een hoogteprofiel, een temperatuurverloop, een grafiek op een scherm. Toch kun je de grafiek van de afgeleide schetsen. Je leest dan geen formule af, maar de vorm: waar stijgt de grafiek, waar daalt ze, waar is ze het steilst?

\voorbeeld{}{De linkse figuur toont het hoogteprofiel van een wandeling: de hoogte $h$ in meter na $x$ km. Schets de grafiek van $h'$. Wat betekent $h'$ hier?\label{vb:h2-schets}}{%
\figuren{\begin{tikzpicture}[x=.6cm,y=.03cm,font=\scriptsize]
  \draw[black!14,very thin,xstep=.6cm,ystep=.75cm] (0,0) grid (6,100);
  \draw[->,grijs] (0,0) -- (6.5,0) node[right,text=black] {$x$};
  \draw[->,grijs] (0,0) -- (0,112) node[above,text=black] {$h$};
  \foreach \t in {1,...,6}{\node[below,inner sep=1.5pt] at (\t,0) {$\t$};}
  \foreach \t/\l in {0/200,50/250,100/300}{\node[left,inner sep=1.5pt] at (0,\t) {$\l$};}
  \kromme{0}{6}{10*((\x)^3/3-3.5*(\x)^2+10*\x)}
  \draw[grijs,dashed] (0,0) -- (.6,60);
\end{tikzpicture}\qquad\qquad
\begin{tikzpicture}[x=.6cm,y=.025cm,font=\scriptsize]
  \draw[black!14,very thin,xstep=.6cm,ystep=.625cm] (0,-25) grid (6,100);
  \draw[->,grijs] (0,0) -- (6.5,0) node[right,text=black] {$x$};
  \draw[->,grijs] (0,-25) -- (0,114) node[above,text=black] {$h'$};
  \foreach \t in {1,...,6}{\node[below,inner sep=1.5pt] at (\t,0) {$\t$};}
  \foreach \t in {-25,50,100}{\node[left,inner sep=1.5pt] at (0,\t) {$\t$};}
  \kromme{0}{6}{10*((\x)^2-7*\x+10)}
\end{tikzpicture}}
\uitleg{Zoek eerst de punten met een horizontale raaklijn, dan het teken van de helling op elk stuk, en dan waar het pad het steilst is.}
\begin{stappenbreed}
h'(2) = 0 \text{ en } h'(5) = 0 & In de top na $2$ km en in het dal na $5$ km is de raaklijn horizontaal: daar snijdt de grafiek van $h'$ de $x$-as. \\[1ex]
h' > 0 \text{ voor } x < 2 \text{ en } x > 5, \quad h' < 0 \text{ ertussen} & Waar het pad stijgt, ligt de grafiek van $h'$ boven de $x$-as; waar het daalt, eronder. \\[1ex]
h'(0) \approx 100 & Bij de start is het pad het steilst: de raaklijn (streepjes) stijgt ongeveer $60$ m op $0{,}6$ km, dus $100$ m per km. \\[1ex]
h' \text{ het kleinst rond } x = 3{,}5 & Halverwege de afdaling daalt het pad het steilst: daar ligt de grafiek van $h'$ het laagst.
\end{stappenbreed}
Met die vier vaststellingen teken je de rechtse grafiek. $h'$ is de helling van het pad, in meter per km: bij de start is dat $100$ m per km, een helling van $10\,\%$.}

\begin{werkwijze}
Zoek de punten waar de grafiek van $f$ een horizontale raaklijn heeft: daar is $f' = 0$. Bepaal op elk stuk ertussen of $f$ stijgt ($f'$ boven de $x$-as) of daalt ($f'$ eronder). Zoek waar $f$ het steilst is: daar ligt $f'$ het verst van de $x$-as. Verbind tot een vloeiende grafiek. Ken je het voorschrift, controleer dan met ICT: GeoGebra tekent de afgeleide als je \texttt{f'(x)} invoert, en op de TI-84 voer je \texttt{Y2 = nDeriv(Y1,X,X)} in (via \texttt{MATH}, \texttt{8:nDeriv}).
\end{werkwijze}

\begin{oefeningen}
\oef{h2-schets}
\end{oefeningen}
\begin{alleenvoor}{2,3}

\subsection{Een breuk met de definitie}

Met de definitie kun je ook de afgeleide functie van een breuk berekenen. Breng dan de teller van het differentiequotiënt eerst op één noemer.

\voorbeeld{}{\alleenvoor{2}{Bereken met de definitie de afgeleide functie van $f(x) = \frac1x$.}\alleenvoor{3}{Bewijs met de definitie dat $\bigl(\frac1x\bigr)' = -\frac{1}{x^2}$.}\label{vb:h2-functie}}{%
\uitleg{De teller is een verschil van twee breuken. Zet het op één noemer: dan verschijnt de factor $h$.}
\begin{stappenbreed}
\frac{1}{x + h} - \frac1x = \frac{x - (x + h)}{x(x + h)} = \frac{-h}{x(x + h)} & Op één noemer. \\[1ex]
\frac{f(x + h) - f(x)}{h} = \frac{-h}{h\,x(x + h)} = \frac{-1}{x(x + h)} & Delen door $h$, en schrappen. \\[1ex]
\lim_{h \to 0} \frac{-1}{x(x + h)} = -\frac{1}{x^2} & Invullen, voor $x \neq 0$.
\end{stappenbreed}}

\begin{oefeningen}
\oef{h2-deffunctie}
\end{oefeningen}

\subsection{Constante en \texorpdfstring{$x$}{x}}

Voor een paar functies berekenen we de afgeleide één keer met de definitie. Daarna gebruiken we het resultaat als regel. We beginnen bij de eenvoudigste.

\begin{theorie}
\begin{stelling}{Afgeleide van een constante en van \texorpdfstring{$x$}{x}}
$\formule{(c)' = 0}$ voor elk getal $c$, \qquad en \qquad $\formule{(x)' = 1}$.
\end{stelling}
\begin{bewijs}
Voor $f(x) = c$ is het differentiequotiënt $\frac{c - c}{h} = 0$, met limiet $0$. Voor $f(x) = x$ is het $\frac{(x + h) - x}{h} = 1$, met limiet $1$.
\end{bewijs}
\end{theorie}

Meetkundig is dat duidelijk: de grafiek van een constante is een horizontale rechte, met richtingscoëfficiënt $0$, en de grafiek van $x$ is de rechte $y = x$, met richtingscoëfficiënt $1$.

\subsection{Machten met een natuurlijke exponent}

Uit $(x^2)' = 2x$ en $(x)' = 1$ kun je al een patroon vermoeden. \alleenvoor{2}{Het geldt voor elke natuurlijke exponent.}\alleenvoor{3}{Het binomium van Newton bewijst het voor elke natuurlijke exponent.}

\begin{theorie}
\begin{stelling}{Afgeleide van een macht}
Voor elk natuurlijk getal $n \geqslant 1$: \quad $\formule{(x^n)' = n\,x^{n - 1}}$.
\end{stelling}
\begin{alleenvoor}{3}
\begin{bewijs}
Met het binomium van Newton is
\bewijsregel{(x + h)^n = x^n + n\,x^{n-1} h + \tbinom{n}{2} x^{n-2} h^2 + \cdots + h^n .}
Alle termen na de tweede bevatten minstens $h^2$. Trek $x^n$ af en deel door $h$:
\bewijsregel{\frac{(x + h)^n - x^n}{h} = n\,x^{n-1} + h \Bigl(\tbinom{n}{2} x^{n-2} + \cdots + h^{n-2}\Bigr).}
Tussen de haakjes staat een veelterm in $h$, met een limiet voor $h \to 0$. Maal $h$ gaat die naar $0$, dus het differentiequotiënt gaat naar $n\,x^{n-1}$.
\end{bewijs}
\end{alleenvoor}
\end{theorie}

In woorden: de exponent komt ervoor, en de nieuwe exponent is één kleiner. Voor $n = 1$ krijg je $(x)' = 1 \cdot x^0 = 1$, wat we al wisten.\alleenvoor{2}{ We bewijzen de regel niet: het bewijs steunt op het binomium van Newton. Je kunt hem wel controleren. Voor $x^3$ vind je met de definitie $3x^2$ (dat is een opdracht verderop), en met ICT zie je dat de hellingsgrafiek van $x^3$ de parabool $y = 3x^2$ is: teken in GeoGebra $f(x) = x^3$ en $f'(x)$.}

\voorbeeld{}{Bereken de afgeleide van $f(x) = x^4$, en de richtingscoëfficiënt van de raaklijn in het punt met $x = -2$.\label{vb:h2-macht}}{%
\uitleg{Pas de regel toe, en vul daarna $-2$ in de afgeleide in, niet in $f$.}
\begin{stappen}
f'(x) = 4x^3 & De exponent $4$ komt ervoor, de nieuwe exponent is $3$. \\[1ex]
f'(-2) = 4 \cdot (-8) = -32 & De raaklijn in $(-2, 16)$ daalt sterk.
\end{stappen}}

\begin{oefeningen}
\oef{h2-machtn}
\end{oefeningen}

\subsection{Machten met een andere exponent}

Wat met $\frac1x = x^{-1}$? Met de definitie vonden we $\bigl(\frac1x\bigr)' = -\frac{1}{x^2} = -1 \cdot x^{-2}$: dat is precies wat de regel geeft voor $n = -1$. De regel blijkt voor elke rationale exponent te gelden.

\begin{theorie}
\begin{stelling}{Afgeleide van een macht, algemeen}
Voor elk rationaal getal $r$: \quad $\formule{(x^r)' = r\,x^{r - 1}}$ \quad (waar $x^r$ en $x^{r-1}$ bestaan).
\end{stelling}
\end{theorie}

\alleenvoor{2}{We bewijzen dit niet. Voor $r = -1$ klopt de regel met wat we met de definitie vonden, en met ICT kun je hem voor andere exponenten controleren: de hellingsgrafiek van $\sqrt x = x^{\frac12}$ is die van $\frac12 x^{-\frac12} = \frac{1}{2\sqrt x}$.}\alleenvoor{3}{We bewijzen dit in stappen. Voor $r = -1$ deden we het al, en voor $r = \frac12$ volgt het hieronder met de definitie. Voor elke negatieve gehele exponent volgt het met de quotiëntregel in sectie~\ref{sec:h2-rekenregels}. Het algemene bewijs, voor $x > 0$, volgt in hoofdstuk 2, als we de afgeleide van $\e^x$ en $\ln x$ kennen.}

\begin{alleenvoor}{3}
\voorbeeld{}{Bewijs met de definitie dat $(\sqrt x)' = \dfrac{1}{2\sqrt x}$ voor $x > 0$.\label{vb:h2-wortel}}{%
\uitleg{In de teller staat een verschil van twee wortels. Vermenigvuldig teller en noemer met de toegevoegde, zoals bij limieten met een wortel.}
\begin{stappenbreed}
\frac{\sqrt{x + h} - \sqrt x}{h} = \frac{(\sqrt{x + h} - \sqrt x)(\sqrt{x + h} + \sqrt x)}{h\,(\sqrt{x + h} + \sqrt x)} & Maal de toegevoegde. \\[1.5ex]
= \frac{(x + h) - x}{h\,(\sqrt{x + h} + \sqrt x)} = \frac{1}{\sqrt{x + h} + \sqrt x} & In de teller staat $a^2 - b^2$; schrap $h$. \\[1.5ex]
\lim_{h \to 0} \frac{1}{\sqrt{x + h} + \sqrt x} = \frac{1}{2\sqrt x} & Invullen.
\end{stappenbreed}
Dat is wat de regel geeft: $\bigl(x^{\frac12}\bigr)' = \frac12 x^{-\frac12} = \frac{1}{2\sqrt x}$.}
\end{alleenvoor}

Met de regel hoef je de definitie niet meer te gebruiken. Schrijf de functie eerst als een macht van $x$.

\voorbeeld{}{Bereken de afgeleide van $f(x) = \dfrac{1}{x^2\sqrt x}$.\label{vb:h2-machtr}}{%
\uitleg{Er staat geen macht, maar een breuk met een wortel. Herschrijf ze als één macht van $x$ met de rekenregels voor machten.}
\begin{stappen}
x^2\sqrt x = x^2 \cdot x^{\frac12} = x^{\frac52} & Machten met hetzelfde grondtal: tel de exponenten op. \\[1ex]
f(x) = x^{-\frac52} & Een macht in de noemer krijgt een negatieve exponent. \\[1ex]
f'(x) = -\tfrac52\,x^{-\frac72} = -\frac{5}{2x^3\sqrt x} & De regel; terugschrijven met $x^{\frac72} = x^3\sqrt x$.
\end{stappen}}

\begin{oefeningen}
\oef{h2-functie}
\end{oefeningen}

%========================================================================================
\section{Rekenregels}\label{sec:h2-rekenregels}

Met de definitie kun je in principe elke afgeleide berekenen, maar het rekenwerk wordt snel zwaar. Daarom \alleenvoor{2}{geven we een paar rekenregels. We bewijzen ze niet, maar we tonen telkens met een figuur of een controle waarom ze kloppen.}\alleenvoor{3}{bewijzen we een paar rekenregels.} Daarmee leid je elke functie af die uit machten is opgebouwd met optellen, vermenigvuldigen en delen.

\subsection{Veelvoud}

Rek je een grafiek verticaal uit met een factor $c$, dan worden alle hellingen ook $c$ keer zo groot.

\begin{theorie}
\begin{stelling}{Veelvoudregel}
Is $f$ afleidbaar en $c$ een getal, dan is \quad $\formule{(c\,f)' = c\,f'}$.
\end{stelling}
\begin{alleenvoor}{3}
\begin{bewijs}
Het differentiequotiënt van $c\,f$ is
\bewijsregel{\frac{c\,f(x + h) - c\,f(x)}{h} = c \cdot \frac{f(x + h) - f(x)}{h} ,}
en een constante factor mag voor de limiet. De limiet is dus $c\,f'(x)$.
\end{bewijs}
\end{alleenvoor}
\end{theorie}

\voorbeeld{}{Bereken de afgeleide van $f(x) = -\dfrac{4}{x}$.\label{vb:h2-veelvoud}}{%
\uitleg{Schrijf de functie als een getal maal een macht. De factor $-4$ blijft staan, en de macht leid je af.}
\begin{stappen}
f(x) = -4\,x^{-1} & Een getal maal een macht. \\[1ex]
f'(x) = -4 \cdot (-1)\,x^{-2} = \frac{4}{x^2} & Veelvoudregel en de regel voor machten.
\end{stappen}}

\begin{oefeningen}
\oef{h2-veelvoud}
\end{oefeningen}

\subsection{Som en verschil}

Tel je twee functies op, dan tellen ook hun veranderingen op.\alleenvoor{2}{ Verandert $f$ ongeveer $f'(x) \cdot h$ en $g$ ongeveer $g'(x) \cdot h$ als $x$ met $h$ toeneemt, dan verandert $f + g$ ongeveer $\bigl(f'(x) + g'(x)\bigr) \cdot h$.}

\begin{theorie}
\begin{stelling}{Somregel}
Zijn $f$ en $g$ afleidbaar, dan is \quad $\formule{(f + g)' = f' + g'}$ \quad en \quad $\formule{(f - g)' = f' - g'}$.
\end{stelling}
\begin{alleenvoor}{3}
\begin{bewijs}
Het differentiequotiënt van $f + g$ is
\bewijsregel{\frac{\bigl(f(x + h) + g(x + h)\bigr) - \bigl(f(x) + g(x)\bigr)}{h} = \frac{f(x + h) - f(x)}{h} + \frac{g(x + h) - g(x)}{h} .}
Beide breuken hebben een limiet, $f'(x)$ en $g'(x)$, dus hun som gaat naar $f'(x) + g'(x)$. Voor het verschil gaat het net zo.
\end{bewijs}
\end{alleenvoor}
\end{theorie}

Met de veelvoud- en de somregel leid je elke veelterm term per term af.

\voorbeeld{}{Bereken de afgeleide van $f(x) = 3x^4 - 2\sqrt x + \dfrac5x$.\label{vb:h2-som}}{%
\uitleg{Schrijf elke term als een getal maal een macht. Dan leid je term per term af, met de som- en de veelvoudregel.}
\begin{stappenbreed}
f(x) = 3x^4 - 2x^{\frac12} + 5x^{-1} & Als machten. \\[1ex]
f'(x) = 3 \cdot 4x^3 - 2 \cdot \tfrac12 x^{-\frac12} + 5 \cdot (-1)\,x^{-2} & Elke term apart, de factoren blijven staan. \\[1ex]
= 12x^3 - \frac{1}{\sqrt x} - \frac{5}{x^2} & Terugschrijven zonder negatieve en gebroken exponenten.
\end{stappenbreed}}

\begin{oefeningen}
\needspace{8\baselineskip}
\oef{h2-som}
\end{oefeningen}

\subsection{Product}

Bij een product ligt de fout voor de hand: de afgeleiden vermenigvuldigen. Maar $x^2 = x \cdot x$ heeft afgeleide $2x$, terwijl $(x)' \cdot (x)' = 1 \cdot 1 = 1$. De afgeleide van een product is dus niet het product van de afgeleiden. De juiste regel heeft twee termen.

\begin{theorie}
\begin{stelling}{Productregel}
Zijn $f$ en $g$ afleidbaar, dan is \quad $\formule{(f\,g)' = f'\,g + f\,g'}$.
\end{stelling}
\begin{alleenvoor}{3}
\begin{bewijs}
Tel $f(x)g(x + h)$ op en trek het af:
\bewijsregel{f(x + h)g(x + h) - f(x)g(x) = \bigl(f(x + h) - f(x)\bigr)g(x + h) + f(x)\bigl(g(x + h) - g(x)\bigr).}
Deel door $h$:
\bewijsregel{\frac{f(x + h)g(x + h) - f(x)g(x)}{h} = \frac{f(x + h) - f(x)}{h} \cdot g(x + h) + f(x) \cdot \frac{g(x + h) - g(x)}{h} .}
Laat $h \to 0$. De breuken gaan naar $f'(x)$ en $g'(x)$. Omdat $g$ afleidbaar is, gaat $g(x + h) - g(x) = \frac{g(x + h) - g(x)}{h} \cdot h$ naar $g'(x) \cdot 0 = 0$, dus $g(x + h) \to g(x)$. Het geheel gaat naar $f'(x)g(x) + f(x)g'(x)$.
\end{bewijs}
\end{alleenvoor}
\end{theorie}

In woorden: de eerste factor afgeleid maal de tweede, plus de eerste maal de tweede afgeleid. Voor $x \cdot x$ geeft dat $1 \cdot x + x \cdot 1 = 2x$, zoals het hoort.\alleenvoor{2}{ Het bewijs laten we weg; het voorbeeld hieronder controleert de regel nog eens.}

\voorbeeld{}{Bereken de afgeleide van $f(x) = (x^2 + 1)(x^3 - x)$ met de productregel. Controleer door eerst uit te werken.\label{vb:h2-product}}{%
\uitleg{Neem $x^2 + 1$ als eerste en $x^3 - x$ als tweede factor. Leid elke factor apart af.}
\begin{stappenbreed}
f'(x) = 2x\,(x^3 - x) + (x^2 + 1)(3x^2 - 1) & Productregel, met $(x^2 + 1)' = 2x$ en $(x^3 - x)' = 3x^2 - 1$. \\[1ex]
= 2x^4 - 2x^2 + 3x^4 + 2x^2 - 1 = 5x^4 - 1 & Uitwerken en samennemen. \\[1ex]
f(x) = x^5 - x \Rightarrow f'(x) = 5x^4 - 1 & Controle: eerst uitwerken, dan afleiden, geeft hetzelfde.
\end{stappenbreed}}

Hier was uitwerken even snel. Maar bij een product met een wortel, of later met \alleenvoor{3}{$\sin x$ of }$\e^x$, kun je niet uitwerken, en dan is de productregel nodig.

\begin{oefeningen}
\oef{h2-product}
\end{oefeningen}

\subsection{Quotiënt}

Voor een quotiënt geldt een regel van dezelfde soort, met een min-teken en een kwadraat in de noemer.

\begin{theorie}
\begin{stelling}{Quotiëntregel}
Zijn $f$ en $g$ afleidbaar en is $g(x) \neq 0$, dan is \quad $\formule{\Bigl(\frac{f}{g}\Bigr)' = \frac{f'\,g - f\,g'}{g^2}}$.
\end{stelling}
\begin{alleenvoor}{3}
\begin{bewijs}
Zoals in het bewijs van de productregel gaat $g(x + h)$ naar $g(x) \neq 0$, dus is ook $g(x + h) \neq 0$ voor kleine $h$. Zet het verschil op één noemer, en tel $f(x)g(x)$ op en trek het af:
\bewijsregel{\frac{f(x + h)}{g(x + h)} - \frac{f(x)}{g(x)} = \frac{\bigl(f(x + h) - f(x)\bigr)g(x) - f(x)\bigl(g(x + h) - g(x)\bigr)}{g(x + h)\,g(x)} .}
Deel door $h$. Dan verschijnen de differentiequotiënten van $f$ en $g$:
\bewijsregel{\frac{1}{h}\Bigl(\frac{f(x + h)}{g(x + h)} - \frac{f(x)}{g(x)}\Bigr) = \Bigl(\frac{f(x + h) - f(x)}{h} \cdot g(x) - f(x) \cdot \frac{g(x + h) - g(x)}{h}\Bigr) \cdot \frac{1}{g(x + h)\,g(x)} .}
Voor $h \to 0$ gaat de factor tussen de haakjes naar $f'(x)g(x) - f(x)g'(x)$, en de laatste factor naar $\frac{1}{g(x)^2}$, want $g(x + h) \to g(x)$.
\end{bewijs}
\end{alleenvoor}
\end{theorie}

In woorden: teller afgeleid maal noemer, min teller maal noemer afgeleid, gedeeld door de noemer in het kwadraat. Let op de volgorde: het min-teken maakt dat de twee termen niet mogen wisselen.\alleenvoor{2}{ Ook deze regel bewijzen we niet. Controle: met teller $1$ en noemer $x$ geeft hij $\frac{0 \cdot x - 1 \cdot 1}{x^2} = -\frac{1}{x^2}$, wat we met de definitie vonden.}

\voorbeeld{}{Bereken de afgeleide van $f(x) = \dfrac{2x + 1}{x - 3}$.\label{vb:h2-quotient}}{%
\uitleg{De teller is $2x + 1$ met afgeleide $2$, de noemer $x - 3$ met afgeleide $1$.}
\begin{stappenbreed}
f'(x) = \frac{2 \cdot (x - 3) - (2x + 1) \cdot 1}{(x - 3)^2} & Quotiëntregel: eerst de afgeleide van de teller. \\[1.5ex]
= \frac{2x - 6 - 2x - 1}{(x - 3)^2} = \frac{-7}{(x - 3)^2} & Werk alleen de teller uit; de noemer laat je staan.
\end{stappenbreed}
De afgeleide is overal negatief: de functie daalt op elk interval van haar domein.}

\begin{oefeningen}
\oef{h2-quotient}
\end{oefeningen}

Met de quotiëntregel kunnen we nu ook de regel voor machten bewijzen voor elke negatieve gehele exponent.

\voorbeeld{}{Toon aan dat $\Bigl(\dfrac{1}{x^n}\Bigr)' = -\dfrac{n}{x^{n + 1}}$ voor elk natuurlijk getal $n \geqslant 1$.\label{vb:h2-negmacht}}{%
\uitleg{Pas de quotiëntregel toe met teller $1$ en noemer $x^n$. Voor de noemer gebruiken we de regel voor natuurlijke exponenten, die al bewezen is.}
\begin{stappenbreed}
\Bigl(\frac{1}{x^n}\Bigr)' = \frac{0 \cdot x^n - 1 \cdot n\,x^{n - 1}}{(x^n)^2} & Quotiëntregel: de afgeleide van de constante teller is $0$. \\[1.5ex]
= \frac{-n\,x^{n - 1}}{x^{2n}} = -\frac{n}{x^{n + 1}} & Schrap $x^{n - 1}$: er blijft $x^{2n - (n - 1)} = x^{n + 1}$ over.
\end{stappenbreed}
Met $r = -n$ is dat $-n\,x^{-n-1} = r\,x^{r - 1}$: de regel voor machten geldt ook voor negatieve gehele exponenten.}

\begin{werkwijze}
Schrijf eerst wat kan als een macht van $x$: dan volstaan de veelvoud- en de somregel. Gebruik de productregel en de quotiëntregel alleen als je niet eenvoudig kunt uitwerken of splitsen. Vereenvoudig op het einde.
\end{werkwijze}

\begin{oefeningen}
\oef{h2-negmacht}

\oef{h2-rekenregels}
\end{oefeningen}

\subsection{Terug naar het probleem}

Met de regel voor machten en de veelvoudregel kunnen we nu de stijgsnelheid in het glas berekenen, zonder de definitie.

\voorbeeld{}{Terug naar het probleem. Je schenkt een kegelvormig cocktailglas vol, $10$ ml per seconde. Na $t$ seconden staat het vocht $h(t) = 3\sqrt[3]{t}$ cm hoog. Hoe snel stijgt het vocht na $8$ seconden? En na $1$ seconde?\label{vb:h2-terug}}{%
\naast{%
\uitleg{De stijgsnelheid op een ogenblik is de afgeleide $h'(t)$. Schrijf de wortel als een macht van $t$: dan volstaan de veelvoudregel en de regel voor machten.}
\begin{stappenvrij}{4.1cm}{4.3cm}
h(t) = 3\,t^{\frac13} & Een getal maal een macht. \\
h'(t) = 3 \cdot \tfrac13\,t^{-\frac23} = \frac{1}{\sqrt[3]{t^2}} & Veelvoudregel en de regel voor machten. \\
h'(8) = \frac{1}{\sqrt[3]{64}} = \frac14 & Na $8$ s stijgt het vocht $0{,}25$ cm per seconde. \\
h'(1) = 1 & Na $1$ s nog $1$ cm per seconde: vier keer zo snel.
\end{stappenvrij}}{%
\begin{tikzpicture}[x=.25cm,y=.45cm,font=\scriptsize]
  \draw[black!14,very thin,xstep=1cm,ystep=.9cm] (0,0) grid (20,8);
  \draw[->,grijs] (0,0) -- (21,0) node[right,text=black] {$t$};
  \draw[->,grijs] (0,0) -- (0,8.8) node[above,text=black] {$h$};
  \foreach \t in {4,8,12,16}{\node[below,inner sep=1.5pt] at (\t,0) {$\t$};}
  \foreach \t in {2,4,6,8}{\node[left,inner sep=1.5pt] at (0,\t) {$\t$};}
  \draw[accent,thick] plot[domain=0:8,samples=60,smooth] ({(\x/3)^3},\x);
  \draw[grijs,thick] (0,4) -- (16,8);
  \vol{(8,6)}
\end{tikzpicture}}
In de figuur is $\frac14$ de richtingscoëfficiënt van de raaklijn in $(8, 6)$. Dat klopt met de tabel in sectie 1, en ook met de vorm van het glas. Op $6$ cm hoogte is het vochtoppervlak $40$ cm$^2$, en $10$ cm$^3$ per seconde verdeeld over $40$ cm$^2$ geeft een laagje van $0{,}25$ cm per seconde. In $t = 0$ is $h$ niet afleidbaar: de grafiek vertrekt verticaal, omdat het glas onderaan een punt heeft.}
\end{alleenvoor}

%========================================================================================
\opdrachtenkop

\begin{alleenvoor}{2,3}
\begin{oefeningen}
\needspace{13\baselineskip}
\oef{h2-o10}

\needspace{9\baselineskip}
\oef{h2-o13}

\needspace{8\baselineskip}
\oef{h2-o2}

\needspace{9\baselineskip}
\oef{h2-o16}

\needspace{10\baselineskip}
\oef{h2-o1}

\needspace{10\baselineskip}
\oef{h2-o14}

\oef{h2-o12}

\oef{h2-o15}

\needspace{8\baselineskip}
\oef{h2-o3}

\needspace{12\baselineskip}
\oef{h2-o4}

\needspace{9\baselineskip}
\oef{h2-o23}

\needspace{14\baselineskip}
\oef{h2-o18}

\oef{h2-o5}

\oef{h2-o20}

\needspace{15\baselineskip}
\oef{h2-o11}

\needspace{12\baselineskip}
\oef{h2-o17}

\oef{h2-o6}

\needspace{12\baselineskip}
\oef{h2-o24}

\needspace{9\baselineskip}
\oef{h2-o19}

\oef{h2-o7}

\oef{h2-o21}

\needspace{9\baselineskip}
\oef{h2-o27}

\oef{h2-o8}

\needspace{10\baselineskip}
\oef{h2-o26}

\oef{h2-o9}

\oef{h2-o29}
\end{oefeningen}
\end{alleenvoor}
\begin{alleenvoor}{1}
\begin{oefeningen}
\needspace{9\baselineskip}
\oef{h2-o13}

\needspace{9\baselineskip}
\oef{h2-o16}

\needspace{8\baselineskip}
\oef{h2-o2}

\needspace{10\baselineskip}
\oef{h2-o14}

\oef{h2-o12}

\oef{h2-o15}

\needspace{14\baselineskip}
\oef{h2-o18}

\oef{h2-o20}

\oef{h2-o6}

\needspace{12\baselineskip}
\oef{h2-o24}

\needspace{9\baselineskip}
\oef{h2-o19}

\needspace{9\baselineskip}
\oef{h2-o30}

\oef{h2-o7}

\oef{h2-o8}

\oef{h2-o29}
\end{oefeningen}

%----------------------------------------------------------------------------------------
% KOPPEL: verbanden met andere leerstof
%----------------------------------------------------------------------------------------
\needspace{20\baselineskip}
\opdrachtsectie{Koppel}

\lkr{Exponentiële en logaritmische functies en Goniometrische functies: kijk na wat al gezien is. Het pakket legt de volgorde niet vast. Beide oefeningen gebruiken ICT.}

\begin{oefeningen}
\oef{h2-koppel-3}

\oef{h2-koppel-4}
\end{oefeningen}
\end{alleenvoor}

\begin{alleenvoor}{2,3}
%----------------------------------------------------------------------------------------
% KOPPEL: verbanden met andere leerstof
%----------------------------------------------------------------------------------------
\needspace{20\baselineskip}
\opdrachtsectie{Koppel}

\lkr{Integralen en Rationale functies: kijk na wat al gezien is. Het pakket legt de volgorde niet vast.}

\begin{oefeningen}
\oef{h2-koppel-1}

\oef{h2-koppel-2}

\end{oefeningen}
\end{alleenvoor}

%----------------------------------------------------------------------------------------
\opdrachtsectie{Test jezelf}

\begin{alleenvoor}{2,3}
\begin{oefeningen}
\oef{h2-t1}

\oef{h2-t9}

\oef{h2-t2}

\oef{h2-t3}

\oef{h2-t4}

\oef{h2-t5}

\oef{h2-t6}

\oef{h2-t7}

\oef{h2-t10}
\end{oefeningen}

\herhaalschema{%
  \fv{oef:h2-t1}{vb:h2-definitie}{oef:h2-defpunt}%
  \fv{oef:h2-t9}{vb:h2-tabel}{oef:h2-tabel}%
  \fv{oef:h2-t2}{vb:h2-raaklijn}{oef:h2-o2}%
  \fv{oef:h2-t3}{vb:h2-som}{oef:h2-som}%
  \fv{oef:h2-t4}{vb:h2-product}{oef:h2-product}%
  \fv{oef:h2-t5}{vb:h2-quotient}{oef:h2-quotient}%
  \fv{oef:h2-t6}{vb:h2-snelheid}{oef:h2-o15}%
  \fv{oef:h2-t7}{vb:h2-functie}{oef:h2-deffunctie}%
  \fv{oef:h2-t10}{vb:h2-schets}{oef:h2-schets}}
\end{alleenvoor}
\begin{alleenvoor}{1}
\begin{oefeningen}
\oef{h2-t1}

\oef{h2-t9}

\oef{h2-t2}

\oef{h2-t11}

\oef{h2-t12}

\oef{h2-t6}

\oef{h2-t10}
\end{oefeningen}

\herhaalschema{%
  \fv{oef:h2-t1}{vb:h2-definitie}{oef:h2-defpunt}%
  \fv{oef:h2-t9}{vb:h2-tabel}{oef:h2-tabel}%
  \fv{oef:h2-t2}{vb:h2-raaklijn}{oef:h2-o2}%
  \fv{oef:h2-t11}{vb:h2-context}{oef:h2-context}%
  \fv{oef:h2-t12}{vb:h2-tweedegraad}{oef:h2-tweedegraad}%
  \fv{oef:h2-t6}{vb:h2-snelheid}{oef:h2-o15}%
  \fv{oef:h2-t10}{vb:h2-schets}{oef:h2-schets}}
\end{alleenvoor}

\ifzelfstudie
\section*{Maak je eigen samenvatting}
Sluit dit boek en schrijf in je schrift: het differentiequotiënt en zijn betekenis (koorde, gemiddelde snelheid), de definitie van $f'(a)$\alleenvoor{2,3}{ en hoe je ze met een tabel benadert, de vergelijking van de raaklijn, wat een afgeleide in een context betekent\alleenvoor{2}{ (ook de elasticiteit)}, de afgeleide functie, de afgeleide van $x^r$, hoe je de grafiek van $f'$ schetst, en de rekenregels\alleenvoor{3}{ met het bewijs van de productregel}}\alleenvoor{1}{ en hoe je ze met een tabel benadert, de vergelijking van de raaklijn, wat een afgeleide in een context betekent (met de eenheid), de afgeleide functie, de afgeleide van een eerste- en een tweedegraadsfunctie, en hoe je de grafiek van $f'$ schetst}. Vergelijk daarna met de kaders.
\fi
